\chapter{Literature Review}
\label{chap:lr}
\chaptermark{Second Chapter Heading}

\section{Refinement Types}
\label{sec:refinement_types_defintion}
Refinement types allow the developer to extend a language's type system
with predicates \cite{jhalaRefinementTypesTutorial2020}. For example, the user can constraint
the regular integer type with the following refinement 
to get the type that represents only natural numbers.
\begin{lstlisting}
    @Refinement(identifier: 'x', predicate: 'x >= 0')
    typedef nat = int;
\end{lstlisting}
With such refinements the developer can easily avoid index out of bounds errors by, for example,
defining the following refinement:
\begin{lstlisting}
    @Refinement(parameters: ['list', 'index'], predicate: 'index <= length(list)')
    T get(List<T> list, int index) => list[index];
\end{lstlisting}
The type checker will ensure that the index provided to the function is smaller than the
length of the list, and if not, will report this operation as unsafe.

\section{Liquid Types and Refinement Type Inference}
\label{sec:liquid_types}
Annotating each variable, each type with refinement information can get
overwhelming and consume more time than the annotations will reduce in error finding.
To address this issue we can implement refinement type inference using \textit{Logically Qualified Data Types} also
known as \textit{Liquid Types} system explored by \textit{Rondon et al.}\cite{rondonLiquidTypes2008}
which will provide the developer with strength of refinement types without additional unnecessary effort.
The inference process involves the following steps:
\begin{enumerate}
    \item Hindley-Milner type inference \cite{damasPrincipalTypeSchemesFunctional}
    \item Constraint generation
    \item Constraint solving using an SMT solver \cite{demouraZ3EfficientSMT2008}
\end{enumerate}
Assume a function \textit{safeHead}, where we want to check safety of \textit{xs[0]} operation:
\begin{lstlisting}
    int safeHead(List<int> xs) => xs.length > 0 ? xs[0] : 0;
\end{lstlisting}
Without type inference we would need to manually specify length property of \textit{xs}
as a refinement, however liquid types can infer the constraint from the program itself. 

\section{Refinement Types in Other Languages}
\label{sec:existing_implementations}
Similar implementations exist for other programming languages. The most similar one to my implementation is
a refinement type system for TypeScript \cite{vekrisRefinementTypesTypeScript2016} by \textit{Vekris et al}.
Their work demonstrates an integration of a refinement type system into an imperative object-oriented language while
retaining support for higher-order functions.
\\One of the main problems the implementation faces is treatment of mutable variables.
Since the value of a variable can change during program execution, refinement predicates can become
invalid after reassignment. To address the issue the authors translate input programs into Static Single Assignment form,
where each variable is assigned only once. This allows their checker to reason about individual variable versions and retain
the information about conditions under which they were assigned.
\\The implementation also uses the Liquid Types \cite{rondonLiquidTypes2008} system which, on par with TypeScript 
being similar to Dart on many accounts \cite{DartlangLanguage2026}, makes this implementation
a good reference point for my project.
\\Other notable examples of such system implementations are \cite[Liquid Haskell]{vazouLiquidHaskellExperienceRefinement2014}, 
\cite[Refinement Types For Ruby]{kazerounianRefinementTypesRuby2018}, and \cite[Flux: Liquid Types for Rust]{lehmannFluxLiquidTypes2023}

\section{SMT Solvers and Z3}
\label{sec:smt_solvers}
Satisfiability Modulo Theories (SMT) is a problem that asks whether a first-order logic
formula is satisfiable with respect to a background theory \cite{barrettChapter33Satisfiability2021}. It extends Boolean Satisfiability (SAT) problem
by adding reasoning such as integers, arrays, bit-verctos and uninterpreted functions. Since SMT solvers
are able of reasoning about relationships between values and operations within theories, they become
particularly useful in software verification, where program properties involve arithmetic constraints, equality connections
and data structures.
\\
de Moura and Bjørner present Z3 \cite{demouraZ3EfficientSMT2008}, an SMT solver developed at Microsoft
Research and designed specifically for applications in software verification. In context of refinement
type checking, an SMT solver can serve as the reasoning engine, responsible for verifying refinement predicates and
generating constraints, making Z3 suitable for proposed refinement type checker for Dart. Using Z3 allows to
delegate logical reasoning to Z3 and focus on the implementation on extracting extracting properties and interpreting
verification results.

% \Blindtext[2]

% \section{Another Section}
% \Blindtext[1]
% \newpage
% With the widespread of computing systems, information processing, and net working, the practice of replacing paper documentation to electronic documentation has become more and more common. Electronic documentation within the workplace has several advantages over the traditional one, such as being easy to share, copy and edit the document. However, these advantages also present a problem. A malefactor, having access to the system, can easily copy and leak the document, without leaving any trace. Such actions are virtually undetectable in most systems, so, the malefactor goes unpunished. In this thesis, we propose one solution to the problem: digital watermarking. 

% Every electronic document within the protected system is marked with an invisible digital watermark, containing information about the user, accessing this particular document. Therefore, in case the protected company discovers the leaked document, they will be able to identify the machine of the malicious person and time when the document was leaked. This will allow inflicting punishment on the malefactor, recovering the costs of the leak, and potentially preventing future ones. 

% \begin{longtable}{c|c}
% \caption[This is the title I want to appear in the List of Tables]{Simulation Parameters} \label{table:secsimulation_params} \\
% \hline
% A & B  \\
% \hline
% \endfirsthead
% \multicolumn{2}{@{}l}{} \\
% \hline
% A & B \\
% \hline
% \endhead
% \hline
%  \textbf{Parameter} & \textbf{Value}\\
%  \hline
%  Number of vehicles & $|\mathcal{V}|$\\
%  \hline
%  Number of RSUs & $|\mathcal{U}|$\\
%  \hline
%  RSU coverage radius & 150 m\\
%  \hline
%  V2V communication radius & 30 m\\
%  \hline
%  Smart vehicle antenna height & 1.5 m\\
%  \hline
%  RSU antenna height & 25 m\\
%  \hline
%  Smart vehicle maximum speed & $v_{max}$ m/s\\
%  \hline
%  Smart vehicle minimum speed & $v_{min}$ m/s\\
%  \hline
%  Common smart vehicle cache capacities & $[50, 100, 150, 200, 250]$ mb\\
%  \hline
%  Common RSU cache capacities & $[5000,1000,1500,2000,2500]$ mb\\
%  \hline
%  Common backhaul rates & $[75, 100, 150]$ mb/s\\
%  \hline
% \end{longtable}

% \begin{figure}[hbt]
% \centering
% \includegraphics[]{figs/images.png}
% \caption{One kernel at $x_s$ (\emph{dotted kernel}) or two kernels at
% $x_i$ and $x_j$ (\textit{left and right}) lead to the same summed estimate
% at $x_s$. This shows a figure consisting of different types of
% lines. Elements of the figure described in the caption should be set in
% italics, in parentheses, as shown in this sample caption.}
% \label{fig:secex}
% \end{figure}

% This description implies several essential properties of the task at hand:
% \begin{enumerate}
%     \item Watermark must contain all necessary information, but still, be placeable and recognizable even on smaller images. The produced watermark must be compact but have the possibility to store enough information. 
%     \item To prevent easy tampering, the watermark must be invisible to the naked eye (and, preferably, to basic image parsing tools). If malefactor does not know about the existence of watermark, they might not even try to remove it and disable it. 
% \end{enumerate}
